\section{Model, Oracle, and Algorithm}
\label{sec:setup}

We write $\mathcal E=\{1,\ldots,m\}$ for the environments and
\[
  \Delta_m:=\left\{w\in\R_+^m:\sum_{e=1}^m w_e=1\right\}.
\]
All vector norms in the mathematical statements are Euclidean.  In Lean they
are represented by the explicit coordinate sums
$\operatorname{sqNorm}(x)=\sum_i x_i^2$ and
$\operatorname{sqDist}(x,y)=\sum_i(x_i-y_i)^2$, rather than by the default
norm on the raw function type $\operatorname{Fin}(p)\to\R$.

\subsection{NegDRO population objectives}

For an environmental pair $(X^e,Y^e)\in\R^p\times\R$, define
\begin{equation}
  R_e(b):=\E\bigl[(Y^e-b^\top X^e)^2\bigr].
  \label{eq:population-risk}
\end{equation}
The original negative-weight set and its simplex reparameterization are
\[
  \mathcal U(\gamma)
  :=\left\{u\in\R^m:\sum_eu_e=1,\ \min_eu_e\ge-\gamma\right\},
  \qquad
  u_e=(1+\gamma m)w_e-\gamma,
\]
where $\gamma\ge0$ and $w\in\Delta_m$.  After division by the positive
constant $1+\gamma m$, set
\begin{align}
  a_\gamma&:=\frac{\gamma}{1+\gamma m},
  \label{eq:a-gamma}\\
  \mathcal L(b,w)&:=\sum_{e=1}^m(w_e-a_\gamma)R_e(b),
  &\Phi(b)&:=\max_{w\in\Delta_m}\mathcal L(b,w),
  \label{eq:unpenalized-negdro}\\
  \mathcal L_\mu(b,w)&:=\mathcal L(b,w)-\mu\|w\|_2^2,
  &\Phi_\mu(b)&:=\max_{w\in\Delta_m}\mathcal L_\mu(b,w).
  \label{eq:penalized-negdro}
\end{align}
Here $\mu\ge0$ is the dual penalty.  The population objectives originate in
the NegDRO formulation of Wang et al.~\cite{wang2026negdro}; the stochastic
recursion studied below is the extension developed in this note.

\subsection{Corrected realized additive-intervention model}
\label{subsec:corrected-reduced-form}

The structural equations are
\begin{equation}
  Y=\beta^{*\top}X+\varepsilon_Y,
  \qquad
  X=B_{YX}Y+B_{XX}X+\varepsilon_X.
  \label{eq:sem}
\end{equation}
The additive-intervention decomposition writes the realized noises as
\begin{equation}
  \varepsilon_Y=\eta_Y,
  \qquad
  \varepsilon_X=\eta_X+\delta_e.
  \label{eq:additive-intervention}
\end{equation}
The final Lean entry point starts from the resulting realized environmental
maps, rather than reconstructing them from acyclicity or a block-inverse
theorem.  Let
\[
  S:=(I-B_{XX})-B_{YX}\beta^{*\top},
  \qquad G^\top S=I.
\]
Direct substitution in \eqref{eq:sem} gives
\begin{equation}
  SX=B_{YX}\varepsilon_Y+\varepsilon_X.
  \label{eq:schur-elimination}
\end{equation}
Consequently the corrected realized maps used by the final theorem are
\begin{align}
  X^e
  &=G^\top\bigl(B_{YX}\eta_Y+\eta_X+\delta_e\bigr),
  \label{eq:corrected-realized-X}\\
  Y^e
  &=\beta^{*\top}X^e+\eta_Y.
  \label{eq:corrected-realized-Y}
\end{align}
The Lean argument named \texttt{GT} represents the matrix denoted here by
$G^\top$.  The inverse identity $G^\top S=I$ explains the structural
elimination above, but it is not a field of the final public assumption
bundle.  That interface starts from the corrected realized maps
\eqref{eq:corrected-realized-X}--\eqref{eq:corrected-realized-Y}; it does not
take $B_{XX}$ or acyclicity as parameters.

The primitive model assumptions require integrability of the coordinate
products needed for the second-moment calculation: $\eta_Y^2$,
$\eta_Y\eta_{X,i}$, $\eta_Y\delta_{e,i}$,
$\eta_{X,i}\eta_{X,j}$, $\eta_{X,i}\delta_{e,j}$,
$\delta_{e,i}\eta_{X,j}$, and $\delta_{e,i}\delta_{e,j}$.  They also require
the systematic/intervention cross moments to vanish, in particular
\begin{equation}
  \E[\eta_Y\delta_e]=0,
  \qquad
  \E[\eta_X\delta_e^\top]=0,
  \label{eq:systematic-intervention-orthogonality}
\end{equation}
with the corresponding coordinate orientations.  These are substantive model
assumptions in the final theorem.

Define the systematic covariate noise
$U:=B_{YX}\eta_Y+\eta_X$ and the common linear coefficient
\begin{equation}
  v_{\rm SEM}:=-G^\top\E[U\eta_Y].
  \label{eq:v-sem}
\end{equation}
Under \eqref{eq:systematic-intervention-orthogonality},
$\E[X^e\eta_Y]=-v_{\rm SEM}$ for every $e$.  With
\begin{equation}
  \sigma_Y^2:=\E[\eta_Y^2],
  \qquad
  \Sigma_e:=\E[X^eX^{e\top}],
  \label{eq:sem-moments}
\end{equation}
the environmental risk therefore has the common expansion
\begin{equation}
  R_e(b)
  =\sigma_Y^2
   +2v_{\rm SEM}^\top(b-\beta^*)
   +(b-\beta^*)^\top\Sigma_e(b-\beta^*).
  \label{eq:risk-expansion}
\end{equation}
The Lean development proves this identity from the actual squared-loss
integral and the realized moment relations; it does not identify two intended
``KL-like'' or ``risk-like'' expressions by name alone.

\begin{remark}[Source sign convention]
Equation~(10) of arXiv:2412.11850v3 forces the positive term
$+G^\top B_{YX}\varepsilon_Y$ in the reduced form
\eqref{eq:corrected-realized-X}.  Appendix~E.1 Eq.~(97) instead prints the
opposite sign in the lower-left inverse block.  The negative cross terms later
printed in the matrix $H$ are also inconsistent with the displayed
positive-block factorization.  The formalization follows the algebra implied
by Eq.~(10) and verifies the corrected positive-sign identity.  It does not
claim to verify the inconsistent printed Eq.~(97).  This is a source
discrepancy, not an inference about the authors' intent.
\end{remark}

\subsection{Heterogeneity witness and feasible geometry}

For $w\in\Delta_m$, define
\begin{equation}
  A(w):=\sum_{e=1}^m\left(w_e-\frac1m\right)\Sigma_e.
  \label{eq:heterogeneity-matrix}
\end{equation}
\begin{assumption}[Fixed spectral witness]
\label{ass:heterogeneity}
There are $w^0\in\Delta_m$ and $\lambda>0$ such that every eigenvalue of the
symmetric matrix $A(w^0)$ is at least $\lambda$.  Equivalently,
\begin{equation}
  A(w^0)\succeq\lambda I_p.
  \label{eq:fixed-witness-condition}
\end{equation}
The fixed vector $w^0$ is used only in the proof; it is not supplied to the
algorithm.
\end{assumption}

\begin{assumption}[Closed convex bounded feasible set]
\label{ass:bounded-domain}
Let $C\subseteq\R^p$ be nonempty, closed, and convex, with
\begin{equation}
  b_{\rm Init},\beta^*\in C,
  \qquad
  \|b\|_2^2\le B^2\quad\text{for every }b\in C,
  \qquad B\ge0.
  \label{eq:bounded-domain}
\end{equation}
Write $P_C$ for the Euclidean nearest-point projection onto $C$.
\end{assumption}

The formalization constructs $P_C$ by transporting coordinates to a Euclidean
Hilbert space.  It proves existence, membership in $C$, uniqueness, the
variational inequality, the strong Pythagorean (Fej\'er) inequality,
nonexpansiveness, continuity, and measurability.  Thus the projection used in
the final recursion is not an abstract Fej\'er interface.  The ball
$C=\{b:\|b\|_2\le B\}$ is one concrete example.  Its classical explicit
formula is
\begin{equation}
  P_C(z)=
  \begin{cases}
    z,&\|z\|_2\le B,\\
    Bz/\|z\|_2,&\|z\|_2>B.
  \end{cases}
  \label{eq:ball-projection}
\end{equation}
The Lean development formalizes the general nearest-point construction used
by the theorem; it does not separately package \eqref{eq:ball-projection} as
a ball-specific theorem.

\subsection{Fresh conditional sampling law}

The stochastic process is defined on a probability space
$(\Omega,\mathcal F,\mathbb P)$ with a pre-round filtration
$(\mathcal F_{t-1})_{t\ge1}$.  At round $t$, the observed environment and data
are $(E_t,X_t,Y_t)$.  Intuitively, $E_t$ is conditionally uniform and the
observation is fresh from its selected environmental law.  The precise final
assumption is the following law-level identity.

\begin{assumption}[Universal fresh conditional law]
\label{ass:online-unbiasedness}
For each used round $t$, environment $e$, and real test function $\varphi$ that
is integrable under the fixed environment law $\nu_e$,
\begin{equation}
  \E\!\left[
    \mathbf 1_{\{E_t=e\}}\varphi(X_t,Y_t)
    \mid\mathcal F_{t-1}
  \right]
  =\frac1m\int\varphi(X^e(\xi),Y^e(\xi))\,d\nu_e(\xi)
  \quad\text{a.e.}
  \label{eq:fresh-conditional-law}
\end{equation}
The selected test on the left is integrable, each $\nu_e$ is a probability
measure, and the observed coordinates have the required a.e. strong
measurability.  In the final additive wrapper, all $\nu_e$ are the common
noise law $\nu$ and $(X^e,Y^e)$ are the realized maps
\eqref{eq:corrected-realized-X}--\eqref{eq:corrected-realized-Y}.
\end{assumption}

Taking $\varphi=1$ gives conditional uniform environment selection.  Taking
$XX$, $XY$, and $Y^2$ monomials gives the selected population moments; fourth
degree tests give the raw moment inequalities used below.  These identities,
together with predictability of the current state, derive conditional
unbiasedness of the actual oracles.  Conditional oracle unbiasedness is not a
separate final assumption.

The filtration is increasing, each conditioning sigma-algebra is contained in
the ambient sigma-algebra, and the round-$t$ primal and dual oracles are
measurable with respect to the post-observation sigma-algebra $\mathcal F_t$.
In Lean, this oracle measurability is assumed for every natural-number index,
whereas the universal fresh law is required only for the used rounds $t\le T$.
The current iterates are consequently $\mathcal F_{t-1}$-measurable.  These
adaptedness and measurability conditions remain explicit assumptions.  The
formalization does not construct a product probability space, a natural
filtration, or an infinite independent stream.

\subsection{Concrete oracles, moments, and integrability}

Assume the environmental coordinates $X_i^e$ and outcomes $Y^e$ are a.e.
strongly measurable.  Separately, assume the environmental fourth moments are
measurable and integrable and
satisfy, for nonnegative $K_X,K_Y$,
\begin{equation}
  \E\|X^e\|_2^4\le K_X,
  \qquad
  \E|Y^e|^4\le K_Y
  \quad\text{for every }e.
  \label{eq:data-fourth-moments}
\end{equation}
These assumptions do not require independence between $X^e$ and $Y^e$.

At the predictable pair $(b_t,w_t)$, the concrete coordinate oracles are
\begin{align}
  [\widehat g_t^b]_i
  &=2m\left(w_{t,E_t}-\frac{\gamma}{1+\gamma m}\right)
       (X_t^\top b_t-Y_t)X_{t,i},
  \label{eq:online-primal-gradient}\\
  [\widehat g_t^{w,\mu}]_j
  &=m(Y_t-X_t^\top b_t)^2\mathbf 1_{\{j=E_t\}}-2\mu w_{t,j}.
  \label{eq:online-stochastic-gradients}
\end{align}
The unpenalized dual oracle is the specialization
$\widehat g_t^{w,0}$.  The universal law gives, a.e.,
\begin{align}
  \E[\widehat g_t^b\mid\mathcal F_{t-1}]
  &=\nabla_b\mathcal L_\mu(b_t,w_t),\\
  \E[\widehat g_t^{w,\mu}\mid\mathcal F_{t-1}]
  &=\nabla_w\mathcal L_\mu(b_t,w_t).
  \label{eq:online-unbiasedness}
\end{align}
The first equality is independent of $\mu$ because the penalty contains no
$b$.

Here the population-gradient notation abbreviates the explicit
finite-coordinate vector used in the formalization.  Lean proves its
coordinate formula, the pairing identities required by the recursion, and
the corresponding radial \texttt{HasDerivAt} statement; it does not expose a
separate full Fr\'echet-gradient theorem.  No stronger differentiability
claim is needed for the convergence inequalities below.

The exact squared-moment constants used by the final theorems are
\begin{align}
  G_b^2
  &:=4m\left(2\sqrt{K_YK_X}+2B^2K_X\right),
  \label{eq:primal-moment-constant}\\
  G_{w,0}^2
  &:=m^2\left(8K_Y+8B^4K_X\right),
  \label{eq:unpenalized-dual-moment-constant}\\
  G_{w,\mu}^2
  &:=2m^2\left(8K_Y+8B^4K_X\right)+8\mu^2.
  \label{eq:penalized-dual-moment-constant}
\end{align}
These symbols already denote squared bounds and are never squared again.
The derivation uses
\begin{align*}
  (X^\top b-Y)^2\|X\|_2^2
  &\le2\|X\|_2^2Y^2+2B^2\|X\|_2^4,\\
  (Y-X^\top b)^4
  &\le8Y^4+8B^4\|X\|_2^4,
\end{align*}
the simplex bounds on $w_t$, and the fresh selected moments.  It yields the
conditional estimates
\begin{align}
  \E[\|\widehat g_t^b\|_2^2\mid\mathcal F_{t-1}]&\le G_b^2,\\
  \E[\|\widehat g_t^{w,0}\|_\infty^2\mid\mathcal F_{t-1}]&\le G_{w,0}^2,\\
  \E[\|\widehat g_t^{w,\mu}\|_\infty^2\mid\mathcal F_{t-1}]&\le G_{w,\mu}^2.
  \label{eq:derived-oracle-bounds}
\end{align}
Here the dual infinity bound is proved coordinatewise; the Lean development
does not identify it definitionally with the default function-space norm.

Oracle integrability is derived, not assumed.  Fresh selected mixed and fourth
moments first give integrability of
$\|\widehat g_t^b\|_2^2$.  Coordinate domination
$|[\widehat g_t^b]_i|^2\le\|\widehat g_t^b\|_2^2$ and the probability-space
implication $L^2\subseteq L^1$ give primal coordinate integrability.  The
dual coordinate-bound process is controlled by
\[
  m(Y_t-X_t^\top b_t)^2+2\mu;
\]
the residual fourth-moment estimate gives integrability of its square, and
the same $L^2\subseteq L^1$ argument gives dual coordinate integrability.
Trajectory feasibility, simplex invariance, and oracle measurability supply
the adaptive parts of these arguments.

\subsection{Coupled projected-SGD/EG recursion and indexing}

We use the paper-style convention $t=1,\ldots,T$.  Lean uses the zero-based
range $r=0,\ldots,T-1$; the note's $(b_t,w_t)$ corresponds to Lean's state at
$r=t-1$.  Initialize
\begin{equation}
  b_1=b_{\rm Init}\in C,
  \qquad
  w_1^{\rm alg}=m^{-1}\mathbf 1.
  \label{eq:algorithm-initialization}
\end{equation}
The superscript distinguishes the algorithmic uniform initialization from the
fixed comparator $w^0$.  For constant positive steps $\eta_b,\eta_w$, define
\begin{align}
  b_{t+1}
  &=P_C\bigl(b_t-\eta_b\widehat g_t^b\bigr),
  \label{eq:common-primal-update}\\
  w_{t+1,j}
  &=\frac{w_{t,j}\exp(\eta_w[\widehat g_t^{w,\mu}]_j)}
          {\sum_k w_{t,k}\exp(\eta_w[\widehat g_t^{w,\mu}]_k)}.
  \label{eq:common-dual-update}
\end{align}
Thus the primal step is descent and the dual step is ascent.  Both oracles use
the same current state and may use the same round sample.  The proof requires
predictability of the displacement, not independence between the primal and
dual stochastic gradients.

The two error summaries used below are
\begin{equation}
  D_T:=\frac1T\sum_{t=1}^T\E\|b_t-\beta^*\|_2^2,
  \qquad
  \bar b_T:=\frac1T\sum_{t=1}^T b_t.
  \label{eq:error-summaries}
\end{equation}
The final assumption interface additionally requires $m,T>0$,
$\eta_b,\eta_w,\lambda>0$, and $\gamma\ge0$.  Penalized results require
$\mu\ge0$.
